% Sección de entrada de usuario que estaba comentada antes
% \\subsection{Entrada de Usuario con \\texttt{std::cin}: Haciendo tus Programas Interactivos}
% \\label{subsec:entrada_usuario}
%
% Hasta ahora, nuestros programas han sido como monólogos: solo hablan ellos (imprimiendo en la consola). Pero para que un programa sea verdaderamente útil e interactivo, necesita poder \"escuchar\" al usuario, es decir, leer datos que el usuario introduce por teclado.
%
% Para esto, C++ nos proporciona \\texttt{std::cin} (abreviatura de \"Character Input\"), que es el compañero de \\texttt{std::cout}. Mientras \\texttt{std::cout} usa el operador \\texttt{<<} (inserción) para enviar datos a la consola, \\texttt{std::cin} usa el operador \\texttt{>>} (extracción) para \\textbf{extraer} datos desde el teclado y guardarlos en una variable.
%
% Veamos un ejemplo sencillo:
%
% \\begin{lstlisting}[language=C++]
% #include <iostream> // Necesario para std::cout y std::cin
% #include <string>   // Necesario para std::string
%
% int main() {
%     // Declaramos una variable para guardar el nombre del usuario
%     std::string nombre; // Usamos std::string para texto, lo veremos en detalle más adelante
%
%     // Pedimos al usuario que introduzca su nombre
%     std::cout << \"Por favor, introduce tu nombre: \";
%
%     // Leemos el nombre introducido por el usuario y lo guardamos en la variable 'nombre'
%     std::cin >> nombre;
%
%     // Saludamos al usuario usando el nombre que ha introducido
%     std::cout << \"Hola, \" << nombre << \"! Bienvenido a C++.\" << std::endl;
%
%     // Pedimos la edad
%     int edad;    std::cout << \"\nIntroduce tu edad: \";
%     std::cin >> edad;
%
%     // Mostramos la edad
%     std::cout << \"Tienes \" << edad << \" años.\" << std::endl;
%
%     return 0;
% }
% \\end{lstlisting}

\\textbf{Consideraciones Importantes con \\texttt{std::cin}:}
\\begin{itemize}
    \\item \\textbf{Tipo de Dato:} \\texttt{std::cin} intentará leer el dato en el formato que espera la variable. Si el usuario introduce texto cuando se espera un número (como en el ejemplo de la edad), \\texttt{std::cin} puede fallar y dejar la variable con un valor incorrecto o cero, y el flujo de entrada en un estado de error. Esto es algo que aprenderemos a manejar con más robustez en capítulos futuros.
    \\item \\textbf{Espacios en Blanco:} Por defecto, \\texttt{std::cin >>} lee hasta el primer espacio en blanco. Esto significa que si el usuario introduce \"Juan Pérez\", solo \"Juan\" será leído en la variable \\texttt{nombre}. Para leer líneas completas con espacios, se utiliza otra función llamada \\texttt{std::getline()}, que exploraremos cuando hablemos más a fondo de cadenas de texto.
\\end{itemize}

Con \\texttt{std::cin}, tus programas pueden ahora interactuar dinámicamente con el usuario, abriendo un mundo de posibilidades para crear aplicaciones más interesantes y útiles.